\documentclass[13pt]{article}
\usepackage{amssymb}
\usepackage{amsmath}

\begin{document}
\section*{\huge{Modo Texto}} 
\normalsize
\begin{enumerate}
\item \textbf{Reproduza o texto seguinte:} 
\begin{center} 
\textbf{Até os Génios se Enganam!} 
\end{center}

Em 1640, \textbf{Pierre de Fermat} conjeturou que os números da forma eram números primos. Mas Fermat foi traído por seus cálculos. Em 1732, \textbf{Euler}, com sua usual habilidade em lidar com números muito grandes, mostrou que  $$2^{2^n} + 1 = 6.700.417*671.$$

Os números da forma $2^{2^n}+1$ ficaram conhecidos como \textbf{números de Fermat}, e os números primos desta forma como \textbf{primos de Fermat}. \\

\item \textbf{Reproduza o texto seguinte:}
Esta é \emph{uma frase} \scriptsize{COM DIVERSOS} \normalsize{\underline{tipos de letra}}.

\item \textbf{Reproduza as listas seguintes:}
\begin{itemize}
\item aaa
\item bbb
\begin{description}
\item[aaa] - é sequência de três a's
\item[bbb] - é sequência de três b's
\end{description}
\item ccc
\begin{description}
\item - ddd
\item - eee
\end{description}
\end{itemize}

\item \textbf{Reproduza a tabela seguinte e a legenda respetiva}

\begin{table}[h!]
\caption{Função}
\centering
\begin{tabular}{r | c | l} \hline
& $f(x)$ Conclusão & \\ \hline
$0<x<12$ & $+$ & Crescente \\
$x=12$ & $0$ & Máximo local \\ 
$x>12$ & - & Decrescente \\ \hline
\end{tabular}
\end{table}

\item \textbf{Reproduza exatamente como vê:}

\begin{table}[h]
\caption{Datas de Nascimento} \begin{center}
\begin{tabular}{| lcr | l |} \hline
Nome & Dia & Mês & Ano \\ \hline
Maria & 2 & 2 & 1987 \\
Joana & 5 & 6 & 1986 \\ \hline
\end{tabular}
\end{center}
\end{table}
\end{enumerate}

\newpage
\section*{\huge{Modo Matemático}}
\begin{enumerate}
\item Expressão:

$\frac{\sqrt{a_1} + \sqrt{2 + a_i}}{4} = \frac{1}{\sqrt{2+a_i}}$

\item Reproduza as seguintes equações:

\begin{equation}\label{xx}
\begin{split}
\frac{x}{2} &   + \frac{x}{3} + \frac{x}{4} = 14 \\ 
&  \frac{x}{2} + \frac{x}{3} = \frac{7+2x}{3}
\end{split}
\end{equation}

\item Somatório:
$$ f(x) = -\frac{1}{2} +  \sum_{i=1}^{n}\left(\frac{x_i - 2}{\sigma}\right)^2$$

\item Integral: \\
$\int_{0}^{1} \sqrt{1-x}dx = \frac{2}{3}(2\sqrt{2}-1)$

\item Combinação: \\ \\
Dois números Binomiais ${n \choose k}$ e ${n \choose j}$, com $n$, $k$, $j \in N$, são iguais se: \\
$k = j$, ou \\
$k + j = n$

\item Sistema de Equações:
$$
\begin{cases}
\frac{1}{x-y} + \frac{1}{x+y} = \frac{3}{4}\\
\frac{1}{x-y} - \frac{1}{x+y} = -\frac{1}{4}
\end{cases}
$$

\item Matriz:\\ \\
$
d= \begin{vmatrix}
2 & 0 & -k & 3 \\
1 & 2 & \sqrt{-2} & -k+2 \\
2 & \mod{2} & -k+3 & \frac{k+4}{2}
\end{vmatrix}
$

\end{enumerate}



\end{document}